\chapter{Tableaux croisés et fréquences conditionnelles}
\section*{Exercice 20.1 -- Compléter des effectifs}
\begin{center}\begin{tabular}{lrrr}\toprule & Casque & Sans casque & Total\\\midrule A & 90&30&120\\ B&40&40&80\\ Total&130&70&200\\\bottomrule\end{tabular}\end{center} Chaque marge est une somme ; le total général se retrouve dans les deux directions.

\section*{Exercice 20.2 -- Choisir le dénominateur}
Dans toute la population : $130/200=65\,\%$. Parmi A : $90/120=75\,\%$. Parmi les porteurs : $90/130=9/13\simeq69{,}23\,\%$. Les deux dernières fractions ont le même numérateur, mais décrivent des populations de référence différentes.

\section*{Exercice 20.3 -- Nominal ou ordinal}
Le canton est qualitatif nominal ; la satisfaction est qualitative ordinale ; le nombre d'enfants est quantitatif discret ; la taille est quantitative continue (mesurée à une précision donnée). Une moyenne est directement pertinente pour les deux derniers. Coder faible, moyenne, forte par $1,2,3$ ne prouve pas que les écarts de satisfaction sont égaux.

\section*{Exercice 20.4 -- Lire des données réelles}
Le canton compte $195\,091$ personnes. La première proportion est $97\,721/195\,091\simeq50{,}09\,\%$. La seconde est $97\,721/339\,161\simeq28{,}81\,\%$. Il s'agit de deux fréquences conditionnelles distinctes. La première décrit la composition du canton ; la seconde la répartition géographique de la population féminine. Ce sont des données datées, pas une estimation de la population présente aujourd'hui.

\section*{Exercice 20.5 -- Filtrer avec ET, OU, NON}
Intersection : $90$. Union : $120+130-90=160$, car les $90$ personnes de l'intersection seraient comptées deux fois. Complément de l'union : $200-160=40$, ce sont les personnes de B sans casque. Les mots ET, OU et NON correspondent à intersection, union et complémentaire.

\section*{Exercice 20.6 -- Une inversion trompeuse}
Non : la proportion de filles qui font du théâtre est $9/18=1/2$. La fraction $9/12=3/4$ signifie que les trois quarts des élèves qui font du théâtre sont des filles. Le tableau complet comporte $9$ filles sans théâtre, $3$ garçons avec théâtre et $9$ garçons sans théâtre.
